\section{Damped Cauchy evolution in the power index}
\label{sec:cauchy}

Fix \(\epsilon>0\). The normalized variables are
\(\delta=D/\ell^2\), \(m=\ell^2M=Z(1,t)\), and
\(\sigma=(t-t_0)/\ell^2\). The two residues contributing to the
forcing are \(12Z(2,t)\) and \(16\delta Z(3,t)\). A separate
contour evaluates \(m\). We treat the two kernels together
\cite{BakerCM05,BakerCMReview}.

Let
\[
 K(p,t)=\sum_{n\in\mathcal N}\frac{c_n(t)\exp(\epsilon(p-n)^2)}{p-n},
 \qquad b\le \Re p\le r ,
\]
where \(\mathcal N\) is a finite set of residue indices. For
\(F\) take \(K=K_\epsilon\), \(\mathcal N=\{2,3\}\),
\(c_2=12,\ c_3=16\delta\) and \((b,r)=(1,4)\). For \(m\) take
\(K=G_\epsilon=\exp(\epsilon(p-1)^2)/(p-1)\) and
\((b,r)=(3/4,3/2)\). If \(\delta=0\), the term at \(3\) is
absent. Set
\[
 \begin{array}{c|c|c}
  K & C_\epsilon & \lambda\\ \hline
  K_\epsilon & (12+16d)\exp(4\epsilon) & 1\\
  G_\epsilon & \exp(\epsilon/4) & 1/4
 \end{array}
\]
so \(|K(q+\ii v,t)|\le C_\epsilon\exp(-\epsilon v^2)/
\sqrt{\lambda^2+v^2}\) for either \(q=b,r\).
Along \(\Im p=\pm Y\) in the strip we have
\begin{equation}
 |K(p,t)|\le C_\epsilon\exp(-\epsilon Y^2)/Y,\qquad Y>0 .
 \label{cau:kernel-bounds}
\end{equation}
Indeed \(\Re (p-n)^2=(\Re p-n)^2-(\Im p)^2\)
and \(|\Re p-n|\ge\lambda\) on either vertical.

For \(V_q(Y)\) upward from \(q-\ii Y\) to \(q+\ii Y\), use
\(V_q\) for its infinite limit and
\[
 \mathcal V_{b,r}(B)=\frac1{2\pi\ii}\left(
                      \int_{V_r} B(p)\dd p-\int_{V_b}B(p)\dd p\right).
\]
The counterclockwise rectangle uses the right upward line minus
the left upward line, the lower edge directed rightward,
and the upper edge directed leftward. Both real edges and
\(\epsilon\) remain fixed.

\begin{theorem}[Complete two-line formulas]\label{thm:damped-cauchy}
Under the hypotheses of Theorem~\ref{thm:complex-time}, one has
\begin{equation}
 \begin{aligned}
 F(t)&=\mathcal V_{1,4}(K_\epsilon Z)
                  =12Z(2,t)+16\delta(t)Z(3,t)
                  =\ell^4(12S_2+16DS_3),\\
 m(t)&=\mathcal V_{3/4,3/2}(G_\epsilon Z)=Z(1,t),\\
 M'&=2M^2-\ell^{-4}F,\qquad
 \frac{\dd m}{\dd\sigma}=2m^2-F,\qquad
 \frac{\dd\delta}{\dd\sigma}=2-4m\delta .
 \end{aligned}
 \label{eq:forcing-contour}
\end{equation}
The line integrals converge absolutely and allow physical time
differentiation under the integral on each compact subchart. The
two-line differences take the same values on chart overlaps.
\end{theorem}

\begin{proof}
The value estimate \eqref{app:value-envelope} gives
\(|Z(p,t)|\le A_b\exp(\Theta|\Im p|)\) on the strip. Multiply by
\eqref{cau:kernel-bounds}. If
\(\Phi_\theta(Y)=\exp(-\epsilon Y^2+\theta Y)\) and
\(\kappa_\theta(Y)=2\epsilon Y-\theta\), the normalized sum of the
two horizontal integrals is bounded by
\((r-b)C_\epsilon A_b\Phi_\Theta(Y)/(\pi Y)\).
For \(Y>\Theta/(2\epsilon)\), the absolute error in the vertical
difference after cutting to \(|\Im p|\le Y\) is bounded by
\begin{equation}
 \frac{2C_\epsilon A_b}{\pi}
            \frac{\Phi_\Theta(Y)}{Y\kappa_\Theta(Y)}.
 \label{eq:vertical-tail}
\end{equation}
This follows from
\(\Phi_\theta(v)\le\Phi_\theta(Y)
\exp(-\kappa_\theta(Y)(v-Y))\) for \(v\ge Y\), bounding
\(1/\sqrt{\lambda^2+v^2}\) by \(1/v\).
Residues on the finite rectangle give the stated values since
the damping equals one at its poles. The horizontal bounds and
vertical tails give the two-line limit. The laws for \(M,m,\delta\)
follow from \eqref{eq:pairlaw} and \eqref{eq:low-orders}, with the
local \(2m^2\) outside the contour.

On a compact chart and fixed real-part strip,
\eqref{app:chart-time} gives
\(|Z_t|\le C_Z(1+|p|)^2\exp(\Theta|\Im p|)\).
The complete physical derivatives of the integrands are
\begin{equation}
 \begin{aligned}
 \partial_t(K_\epsilon Z)&=K_\epsilon Z_t+
          \frac{16(D'/\ell^2)\exp(\epsilon(p-3)^2)}{p-3}Z,\\
 \partial_t(G_\epsilon Z)&=G_\epsilon Z_t .
 \end{aligned}
 \label{cau:kernel-time}
\end{equation}
The kernel bounds and Gaussian moments justify the interchange
and the continuity of the derivative. The exact bounds on
differentiated horizontal edges and vertical tails are in
Appendix~\ref{app:cauchy-errors}.

On an overlap the finite difference \eqref{cp:transition} and
its time derivative vanish at the residue positions.
Consequently \(K(Z_B-Z_A)\) has removable poles, and its
physical derivative has removable poles as well. Their complete
rectangle integrals vanish. Each bound used for \(KZ\), or
its derivative, applies to the difference with the sum of the two
chart constants. Passing to the vertical limits proves equality
of the full right-minus-left values and their derivatives. In
particular the integer evolution glues through the compact time
cover.
\end{proof}

\subsection{Divisor cutoff and independent spatial boundary cutoff}

For integer \(N>4T+8\), let \(Z_N\) retain \(|\rho_j|\le N\)
with all ties and multiplicities. The estimates used to prove the
theorem apply also to \(Z_N\). The strip-wide bound in
Appendix~\ref{app:cauchy-errors} gives absolute convergence in
the line integrals as \(N\to\infty\). The sharper error is the
signed residue of \(Z-Z_N\). If \(F_N=\mathcal V_{1,4}(K_\epsilon Z_N)\)
then
\begin{equation}
 \begin{aligned}
 F-F_N&=12[Z(2)-Z_N(2)]+16\delta[Z(3)-Z_N(3)],\\
 |F-F_N|&\le12R_2(N)+16d R_3(N),\\
 R_n(N)&=2\Cu(\ell^2q_0)^n E_1(2n-1)3^{2n-1}
                               N^{1-2n}\Lambda(N).
 \end{aligned}
 \label{cau:residue-error}
\end{equation}
Here \(q_0=16/15\), \(\Lambda(x)=\log(3+2T+|x|)\), and \(E_1\)
is the explicit dyadic constant of \eqref{app:dyadic}.
This follows by applying the residue argument to the whole
difference before using \eqref{app:divisor-tail} at integer \(n\).
The rates are \(\ell^4N^{-3}\log N\) and \(d\ell^6N^{-5}\log N\)
at fixed \(T\). The separate \(m\) evaluation has error at most
\(R_1(N)\). For truncated vertical integrals add
\eqref{eq:vertical-tail} to these allowances.

For a spatial cutoff \(X\ge7\), truncate only the horizontal rays
of \eqref{cp:spatial-boundary} at \(|\Re w|=X\), retaining the
finite connectors and circles. Its derivative uses
\((f_p)_t h+f_p h_t\) from \eqref{eq:complex-time}. The
integrated-by-parts form on the cut rays carries the primitive
endpoint terms. The independent derivative errors
\eqref{app:spatial-tails} and \eqref{app:weighted-spatial-tail}
pay this cutoff. They are uniform at fixed height on compact
time subcharts. The complete value limit may take \(N\) first
on each finite rectangle and then \(Y\), or use the absolute
bounds with \(N,Y\) tending to infinity together at fixed
\(\epsilon,T\). For the time derivative one may take \(X\) first
on finite rectangles and then \(Y\), or take \(X\to\infty\) in
the full vertical difference. These distinct limits use the
derivative of the spatial boundary formula and the values of
the sharp divisor cutoff.

\subsection{Folded integrals and branch corrections}

For either kernel and real physical time, \(K^\sharp=K\).
On the compatible charts \eqref{cp:schwarz} the weighted
integrand \(B=KZ\) has \(B^\sharp=B\) and
\begin{equation}
 \mathcal V_{b,r}(B)=\frac1\pi\int_0^\infty
              \Re\big[B(r+\ii v)-B(b+\ii v)\big]\dd v .
 \label{cau:fold}
\end{equation}
On finite rectangles the two horizontal integrals still enter.
To fold a full rectangle, split it into the upper arc from \(r\)
to \(b\) and the lower arc from \(b\) to \(r\), including their
horizontal edges. Reflection sends the upper arc to the reversed
lower arc, while conjugation of \(1/(2\pi\ii)\) changes its sign.
Thus for a Schwarz integrand the rectangle integral is twice the
real part of the normalized upper-arc integral. In the defective
case add the normalized lower-arc integral of
\(B_+-B_+^\sharp\).
For the \(+\pi\) slice prescription, the weighted defect of
\(B_+=KZ_+\) is
\[
 \mathcal D_B=B_+-B_+^\sharp
   =-2\ii K(p,t)\ell^{2p}\sin(\pi p)\sum_{j:Q_j<0}|Q_j|^{-p}.
\]
In fact \(B_+(q-\ii v)=\overline{B_+(q+\ii v)}+\mathcal D_B(q-\ii v)\).
For an upward line this gives
\begin{equation}
 \frac1{2\pi\ii}\int_{V_q(Y)} B_+(p)\dd p
 =\frac1{\pi}\int_0^Y\Re B_+(q+\ii v)\dd v
  +\frac1{2\pi}\int_0^Y\mathcal D_B(q-\ii v)\dd v .
 \label{cau:weighted-defect}
\end{equation}
The two-sided contour takes the right-minus-left difference
of both terms. The removable singularities of \(\mathcal D_B\)
at the integer residues give zero integral of the defect on a
complete rectangle. Its horizontal limit is zero by the same
Gaussian bound (the sum of the bounds for \(B_+\) and \(B_+^\sharp\)
suffices). Thus the \emph{full} defect line integrals agree on
the two lines. The lower-half defect integrals in
\eqref{cau:weighted-defect} are kept in a folded formula.

